
\documentclass[manuscript]{aastex631}
\usepackage{amsmath}

\newcommand{\Nsub}{\ensuremath{N_\mathrm{sub}}}
\newcommand{\Nchan}{\ensuremath{N_\mathrm{chan}}}
\newcommand{\Nharm}{\ensuremath{N_\mathrm{harm}}}
\newcommand{\Eqn}[1]{Equation~\ref{eqn:#1}}

%% The default is a single spaced, 10 point font, single spaced article.
%% There are 5 other style options available via an optional argument. They
%% can be invoked like this:
%%
%% \documentclass[arguments]{aastex631}
%% 
%% where the layout options are:
%%
%%  twocolumn   : two text columns, 10 point font, single spaced article.
%%                This is the most compact and represent the final published
%%                derived PDF copy of the accepted manuscript from the publisher
%%  manuscript  : one text column, 12 point font, double spaced article.
%%  preprint    : one text column, 12 point font, single spaced article.  
%%  preprint2   : two text columns, 12 point font, single spaced article.
%%  modern      : a stylish, single text column, 12 point font, article with
%% 		  wider left and right margins. This uses the Daniel
%% 		  Foreman-Mackey and David Hogg design.
%%  RNAAS       : Supresses an abstract. Originally for RNAAS manuscripts 
%%                but now that abstracts are required this is obsolete for
%%                AAS Journals. Authors might need it for other reasons. DO NOT
%%                use \begin{abstract} and \end{abstract} with this style.
%%
%% Note that you can submit to the AAS Journals in any of these 6 styles.
%%
%% There are other optional arguments one can invoke to allow other stylistic
%% actions. The available options are:
%%
%%   astrosymb    : Loads Astrosymb font and define \astrocommands. 
%%   tighten      : Makes baselineskip slightly smaller, only works with 
%%                  the twocolumn substyle.
%%   times        : uses times font instead of the default
%%   linenumbers  : turn on lineno package.
%%   trackchanges : required to see the revision mark up and print its output
%%   longauthor   : Do not use the more compressed footnote style (default) for 
%%                  the author/collaboration/affiliations. Instead print all
%%                  affiliation information after each name. Creates a much 
%%                  longer author list but may be desirable for short 
%%                  author papers.
%% twocolappendix : make 2 column appendix.
%%   anonymous    : Do not show the authors, affiliations and acknowledgments 
%%                  for dual anonymous review.
%%
%% these can be used in any combination, e.g.
%%
%% \documentclass[twocolumn,linenumbers,trackchanges]{aastex631}
%%
%% AASTeX v6.* now includes \hyperref support. While we have built in specific
%% defaults into the classfile you can manually override them with the
%% \hypersetup command. For example,
%%
%% \hypersetup{linkcolor=red,citecolor=green,filecolor=cyan,urlcolor=magenta}
%%
%% will change the color of the internal links to red, the links to the
%% bibliography to green, the file links to cyan, and the external links to
%% magenta. Additional information on \hyperref options can be found here:
%% https://www.tug.org/applications/hyperref/manual.html#x1-40003
%%
%% Note that in v6.3 "bookmarks" has been changed to "true" in hyperref
%% to improve the accessibility of the compiled pdf file.
%%
%% If you want to create your own macros, you can do so
%% using \newcommand. Your macros should appear before
%% the \begin{document} command.
%%
\newcommand{\vdag}{(v)^\dagger}
\newcommand\aastex{AAS\TeX}
\newcommand\latex{La\TeX}

%% Reintroduced the \received and \accepted commands from AASTeX v5.2
%\received{March 1, 2021}
%\revised{April 1, 2021}
%\accepted{\today}

%% Command to document which AAS Journal the manuscript was submitted to.
%% Adds "Submitted to " the argument.
%\submitjournal{PSJ}

%% For manuscript that include authors in collaborations, AASTeX v6.31
%% builds on the \collaboration command to allow greater freedom to 
%% keep the traditional author+affiliation information but only show
%% subsets. The \collaboration command now must appear AFTER the group
%% of authors in the collaboration and it takes TWO arguments. The last
%% is still the collaboration identifier. The text given in this
%% argument is what will be shown in the manuscript. The first argument
%% is the number of author above the \collaboration command to show with
%% the collaboration text. If there are authors that are not part of any
%% collaboration the \nocollaboration command is used. This command takes
%% one argument which is also the number of authors above to show. A
%% dashed line is shown to indicate no collaboration. This example manuscript
%% shows how these commands work to display specific set of authors 
%% on the front page.
%%
%% For manuscript without any need to use \collaboration the 
%% \AuthorCollaborationLimit command from v6.2 can still be used to 
%% show a subset of authors.
%
%\AuthorCollaborationLimit=2
%
%% will only show Schwarz & Muench on the front page of the manuscript
%% (assuming the \collaboration and \nocollaboration commands are
%% commented out).
%%
%% Note that all of the author will be shown in the published article.
%% This feature is meant to be used prior to acceptance to make the
%% front end of a long author article more manageable. Please do not use
%% this functionality for manuscripts with less than 20 authors. Conversely,
%% please do use this when the number of authors exceeds 40.
%%
%% Use \allauthors at the manuscript end to show the full author list.
%% This command should only be used with \AuthorCollaborationLimit is used.

%% The following command can be used to set the latex table counters.  It
%% is needed in this document because it uses a mix of latex tabular and
%% AASTeX deluxetables.  In general it should not be needed.
%\setcounter{table}{1}

%%%%%%%%%%%%%%%%%%%%%%%%%%%%%%%%%%%%%%%%%%%%%%%%%%%%%%%%%%%%%%%%%%%%%%%%%%%%%%%%
%%
%% The following section outlines numerous optional output that
%% can be displayed in the front matter or as running meta-data.
%%
%% If you wish, you may supply running head information, although
%% this information may be modified by the editorial offices.
%\shorttitle{AASTeX v6.3.1 Sample article}
%\shortauthors{Schwarz et al.}
%%
%% You can add a light gray and diagonal water-mark to the first page 
%% with this command:
%% \watermark{text}
%% where "text", e.g. DRAFT, is the text to appear.  If the text is 
%% long you can control the water-mark size with:
%% \setwatermarkfontsize{dimension}
%% where dimension is any recognized LaTeX dimension, e.g. pt, in, etc.
%%
%%%%%%%%%%%%%%%%%%%%%%%%%%%%%%%%%%%%%%%%%%%%%%%%%%%%%%%%%%%%%%%%%%%%%%%%%%%%%%%%
%\graphicspath{{./}{figures/}}
%% This is the end of the preamble.  Indicate the beginning of the
%% manuscript itself with \begin{document}.

\begin{document}

%\title{Modelling the Cyclic Spectrum}

%\author[0000-0003-2519-7375]{W. van Straten}
%\author{S. Os{\l}owski}
%\author{M. Walker}
%\affil{Manly Astrophysics, 15/41-42 East Esplanade, Manly, NSW 2095, Australia}
%\email{vanstraten.willem@gmail.com}

%% Note that the \and command from previous versions of AASTeX is now
%% depreciated in this version as it is no longer necessary. AASTeX 
%% automatically takes care of all commas and "and"s between authors names.

%% AASTeX 6.31 has the new \collaboration and \nocollaboration commands to
%% provide the collaboration status of a group of authors. These commands 
%% can be used either before or after the list of corresponding authors. The
%% argument for \collaboration is the collaboration identifier. Authors are
%% encouraged to surround collaboration identifiers with ()s. The 
%% \nocollaboration command takes no argument and exists to indicate that
%% the nearby authors are not part of surrounding collaborations.

%% Mark off the abstract in the ``abstract'' environment. 

%\begin{abstract}

%The appendix re-derives the gradient of the merit function, including the partial derivatives of the maximum likelihood estimate of the intrinsic pulse profile.

% The AAS Journals have a 250 word limit for the  abstract.

% \end{abstract}

%% Keywords should appear after the \end{abstract} command. 
%% The AAS Journals now uses Unified Astronomy Thesaurus concepts:
%% https://astrothesaurus.org
%% You will be asked to selected these concepts during the submission process
%% but this old "keyword" functionality is maintained in case authors want
%% to include these concepts in their preprints.
%\keywords{pulsars}

%% From the front matter, we move on to the body of the paper.
%% Sections are demarcated by \section and \subsection, respectively.
%% Observe the use of the LaTeX \label
%% command after the \subsection to give a symbolic KEY to the
%% subsection for cross-referencing in a \ref command.
%% You can use LaTeX's \ref and \label commands to keep track of
%% cross-references to sections, equations, tables, and figures.
%% That way, if you change the order of any elements, LaTeX will
%% automatically renumber them.
%%
%% We recommend that authors also use the natbib \citep
%% and \citet commands to identify citations.  The citations are
%% tied to the reference list via symbolic KEYs. The KEY corresponds
%% to the KEY in the \bibitem in the reference list below. 

%\section{Introduction} \label{sec:intro}

%Stuff

\appendix

%%%%%%%%%%%%%%%%%%%%%%%%%%%%%%%%%%%%%%%%%%%%%%%%%%%%%%%%%%%%%%%%%%%%%%%%%%%%%%%%%%%%
%%%%%%%%%%%%%%%%%%%%%%%%%%%%%%%%%%%%%%%%%%%%%%%%%%%%%%%%%%%%%%%%%%%%%%%%%%%%%%%%%%%%

%%  Gradient of the Merit Function

%%%%%%%%%%%%%%%%%%%%%%%%%%%%%%%%%%%%%%%%%%%%%%%%%%%%%%%%%%%%%%%%%%%%%%%%%%%%%%%%%%%%
%%%%%%%%%%%%%%%%%%%%%%%%%%%%%%%%%%%%%%%%%%%%%%%%%%%%%%%%%%%%%%%%%%%%%%%%%%%%%%%%%%%%

\section{Gradient of the Merit Function} \label{app:gradient}

The model cyclic spectrum measured as a function of pulse harmonic $\alpha$ and radio frequency $\nu$ at time $t$,
%
\begin{equation}
    S^\prime_\mu(\alpha,\nu;t) = H(\nu+\alpha/2;t) H^*(\nu-\alpha/2;t) S_\mu(\alpha) g(t),
\end{equation}
%
where $H(\nu;t)$ is the frequency response of the turbulent interstellar medium along the Earth-pulsar line of sight at time $t$ and $S_\mu(\alpha)$ is the temporal average of the least-squares (maximum likelihood) estimates of the intrinsic pulse profile.
%
Given the observed cyclic spectra, $D_\mu(\alpha,\nu;t)$, the residual difference between the model and observed cyclic spectrum measured at time $t$
%
\begin{equation}
    R_\mu(\alpha,\nu;t) \equiv S^\prime_\mu(\alpha,\nu;t) - D_\mu(\alpha,\nu;t)
\end{equation}
%
is used to define the merit function
%
\begin{equation}
    M = \sum_{t,\alpha,\nu,\mu} R^*_\mu(\alpha,\nu;t) R_\mu(\alpha,\nu;t)
    \label{eqn:merit_function}
\end{equation}
%

%%%%%%%%%%%%%%%%%%%%%%%%%%%%%%%%%%%%%%%%%
%%%%%%%%%%%%%%%%%%%%%%%%%%%%%%%%%%%%%%%%%

%%  Gain variations

%%%%%%%%%%%%%%%%%%%%%%%%%%%%%%%%%%%%%%%%%
%%%%%%%%%%%%%%%%%%%%%%%%%%%%%%%%%%%%%%%%%

\subsection{Gain variations}
\label{sec:gain_variations}

\noindent
Where $M$ is minimized, the partial derivative of $M$ with respect to $g(t)$ is zero; i.e.
\begin{equation}
    \frac{dM}{dg(t)}
    = \sum_{\alpha,\nu,\mu} 2\mathrm{Re} \left[
    R_\mu(\alpha,\nu;t) H^*(\nu+\alpha/2;t) H(\nu-\alpha/2;t) S^*_\mu(\alpha) \right] = 0;
\label{eqn:dM_dg}
\end{equation}
which is solved by
\begin{equation}
g(t) = \frac
{\sum_{\alpha,\nu,\mu} \mathrm{Re} \left[D_\mu(\alpha,\nu;t) H^*(\nu+\alpha/2;t) H(\nu-\alpha/2;t) S^*_\mu(\alpha) \right]}
{\sum_{\alpha,\nu,\mu} |H(\nu+\alpha/2;t)|^2 |H(\nu-\alpha/2;t)|^2 |S^*_\mu(\alpha)|^2 }
\end{equation}


%%%%%%%%%%%%%%%%%%%%%%%%%%%%%%%%%%%%%%%%%
%%%%%%%%%%%%%%%%%%%%%%%%%%%%%%%%%%%%%%%%%

%%  Degenerate phase

%%%%%%%%%%%%%%%%%%%%%%%%%%%%%%%%%%%%%%%%%
%%%%%%%%%%%%%%%%%%%%%%%%%%%%%%%%%%%%%%%%%

\subsection{Degenerate phase}
\label{sec:degenerate_phase}

\noindent

Consider $H'(\nu;t) = H'(\nu;t) z(t)$, where $z(t)=\exp(i\phi_t)$ represents the degenerate phase, and

\begin{equation}
    S''_\mu(\alpha,\nu;t) = H'(\nu+\alpha/2;t) H^{\prime *}(\nu-\alpha/2;t) S_\mu(\alpha) g(t) = S'_\mu(\alpha,\nu;t) zz^*,
\end{equation}

Noting that $zz^*=1$, the partial derivatives of $M'$ with respect to $\phi_t$ must be zero; i.e.
\begin{equation}
    \frac{dM'}{d\phi_t} = \frac{dM'}{dz^*}\frac{dz^*}{d\phi_t}
    = \sum_{\alpha,\nu,\mu} 2\mathrm{Re} \left[
    R'_\mu(\alpha,\nu;t) S^{\prime *}_\mu(\alpha,\nu;t) \right] = 0;
\label{eqn:dM_dphi}
\end{equation}
that is
\begin{equation}
\sum_{\alpha,\nu,\mu} \mathrm{Re} \left[D_\mu(\alpha,\nu;t) S^{\prime *}_\mu(\alpha,\nu;t) \right] =
\sum_{\alpha,\nu,\mu} |S'_\mu(\alpha)|^2
\end{equation}

%%%%%%%%%%%%%%%%%%%%%%%%%%%%%%%%%%%%%%%%%
%%%%%%%%%%%%%%%%%%%%%%%%%%%%%%%%%%%%%%%%%

%%  Intrinsic cyclic spectrum

%%%%%%%%%%%%%%%%%%%%%%%%%%%%%%%%%%%%%%%%%
%%%%%%%%%%%%%%%%%%%%%%%%%%%%%%%%%%%%%%%%%

\subsection{Intrinsic cyclic spectrum}
\label{sec:intrinsic_cyclic_spectrum}

\noindent
Where $M$ is minimized, the partial derivative of $M$ with respect to $S_\mu^*(\alpha)$ is zero; i.e.
\begin{equation}
    \frac{dM}{dS_\mu^*(\alpha)}
    = \sum_{t,\nu}
    R_\mu(\alpha,\nu;t) H^*(\nu+\alpha/2;t) H(\nu-\alpha/2;t) g(t) = 0;
\label{eqn:dM_dSx}
\end{equation}
which is solved by
\begin{equation}
S_\mu(\alpha) = \frac
{\sum_{t,\nu} D(\alpha,\nu;t) H^*(\nu+\alpha/2;t) H(\nu-\alpha/2;t) g(t)}
{\sum_{t,\nu} |H(\nu+\alpha/2;t)|^2 |H(\nu-\alpha/2;t)|^2 g^2(t)}
\label{eqn:ml_intrinsic}
\end{equation}

%%%%%%%%%%%%%%%%%%%%%%%%%%%%%%%%%%%%%%%%%
%%%%%%%%%%%%%%%%%%%%%%%%%%%%%%%%%%%%%%%%%

%%  Dynamic impulse response function

%%%%%%%%%%%%%%%%%%%%%%%%%%%%%%%%%%%%%%%%%
%%%%%%%%%%%%%%%%%%%%%%%%%%%%%%%%%%%%%%%%%

\subsection{Dynamic impulse response function}

Now consider the impulse response function at time, $t$, defined by the inverse DFT of $H(\nu;t)$; i.e.
\begin{equation}
    h(\tau;t) = \Nchan^{-1/2} \sum_k H(\nu_k; t) \exp(2\pi i \tau \nu_k)
    \label{eqn:inverse_dft}
\end{equation}
and the inverse relation defined by 
\begin{equation}
    H(\nu;t) = \Nchan^{-1/2} \sum_j h(\tau_j; t) \exp(- 2\pi i \tau_j \nu)
    \label{eqn:forward_dft}
\end{equation}
%
Note that these definitions differ from those presented in equations (31) and (32) of \cite{wdv13},
which normalize only the forward DFT by \Nchan.  Equations~\ref{eqn:inverse_dft} and~\ref{eqn:forward_dft}
conserve energy, such that Parseval's theorem applies; i.e.
\begin{equation}
\sum_{\tau,t} |h(\tau;t)|^2 = \sum_{\nu,t} |H(\nu;t)|^2
\end{equation}

As described in more detail in Appendix A of \cite{ow23}, use these definitions to compute the partial derivative of $M$ with respect to $h^*(\tau_j;t_l)$
\begin{equation}
    \frac{dM}{dh^*(\tau_j;t_l)}
    = \sum_{\alpha,\nu}
    \frac{dS_\mu^{\prime*}(\alpha,\nu;t_l)}{dh^*(\tau_j;t_l)} R(\alpha,\nu;t_l) +
    R^*(\alpha,\nu;t_l) \frac{dS^\prime_\mu(\alpha,\nu;t_l)}{dh^*(\tau_j;t_l)} 
\label{eqn:dM_dh}
\end{equation}
%
where
\begin{equation}
\frac{dS_\mu^{\prime*}(\alpha,\nu;t_l)}{dh^*(\tau_j;t_l)}
    = H(\nu-\alpha/2;t_l) \left( \frac{dH^*(\nu+\alpha/2;t_l)}{dh^*(\tau_j;t_l)} S_\mu^*(\alpha)
                                + H^*(\nu+\alpha/2;t_l) \frac{dS_\mu^*(\alpha)}{dh^*(\tau_j;t_l)} \right) g(t_l)
\label{eqn:dSconj_dh_degenerate}
\end{equation}
%
and
\begin{equation}
\frac{dS^\prime_\mu(\alpha,\nu;t_l)}{dh^*(\tau_j;t_l)}
    = H(\nu+\alpha/2;t_l) \left( \frac{dH^*(\nu-\alpha/2;t_l)}{dh^*(\tau_j;t_l)} S_\mu(\alpha)
                                + H^*(\nu-\alpha/2;t_l) \frac{dS_\mu(\alpha)}{dh^*(\tau_j;t_l)} \right) g(t_l).
\label{eqn:dS_dh_degenerate}
\end{equation}
%

As discussed in Section 2.1 of \cite{wdv13}, it is not possible to minimize $M$ for both $h(\tau,t)$
and $S_\mu^*(\alpha)$ owing to degeneracies that arise in the definition of amplitude, phase, and shift.
Therefore, the second terms in the above equations are dropped, leaving

\begin{equation}
\frac{dS_\mu^{\prime*}(\alpha,\nu;t_l)}{dh^*(\tau_j;t_l)}
    = H(\nu-\alpha/2;t_l) \frac{dH^*(\nu+\alpha/2;t_l)}{dh^*(\tau_j;t_l)} S_\mu^*(\alpha) g(t_l)
\label{eqn:dSconj_dh}
\end{equation}
%
and
\begin{equation}
\frac{dS^\prime_\mu(\alpha,\nu;t_l)}{dh^*(\tau_j;t_l)}
    = H(\nu+\alpha/2;t_l) \frac{dH^*(\nu-\alpha/2;t_l)}{dh^*(\tau_j;t_l)} S_\mu(\alpha) g(t_l).
\label{eqn:dS_dh}
\end{equation}


Noting that
\begin{equation}
    \frac{dH^*(\nu+\alpha/2;t_l)}{dh^*(\tau_j;t_l)} = \Nchan^{-1/2} \exp(2\pi i \tau_j \nu) \exp(\pi i \tau_j \alpha),
\end{equation}
%
substitution of \Eqn{dSconj_dh} into the first term of \Eqn{dM_dh} yields
\begin{equation}
\sum_\alpha \left[ \Nchan^{-1/2} \sum_\nu R(\alpha,\nu;t_l) H(\nu-\alpha/2;t_l) \exp(2\pi i \tau_j \nu) \right] S_\mu^*(\alpha) \exp(\pi i \tau_j \alpha) g(t_l),
\end{equation}
%
which is equivalent to Equation (37) of Walker, Demorest \& van Straten (2013), apart from a factor of 4.
The part of the equation inside the brackets is the inverse DFT of
%
$R(\alpha,\nu;t_l)H(\nu-\alpha/2;t_l)$ 
%
evaluated at $\tau_j$.
%
Likewise, substitution of \Eqn{dS_dh} into the second term of \Eqn{dM_dh} yields
\begin{equation}
\sum_\alpha \left[ \Nchan^{-1/2} \sum_\nu R^*(\alpha,\nu;t_l) H(\nu+\alpha/2;t_l) \exp(2\pi i \tau_j \nu) \right] S_\mu(\alpha) \exp(-\pi i \tau_j \alpha) g(t_l),
\end{equation}
%
which is equivalent to the additional term that was added to the Cyclic-Modeling and pycyc packages.

%%%%%%%%%%%%%%%%%%%%%%%%%%%%%%%%%%%%%%%%%
%%%%%%%%%%%%%%%%%%%%%%%%%%%%%%%%%%%%%%%%%

%%  Doppler-delay wavefield

%%%%%%%%%%%%%%%%%%%%%%%%%%%%%%%%%%%%%%%%%
%%%%%%%%%%%%%%%%%%%%%%%%%%%%%%%%%%%%%%%%%

\subsection{Doppler-delay wavefield}

Now consider the wavefield as a function of delay $\tau$
and Doppler shift $\omega$  defined by 
\begin{equation}
    h(\tau;\omega) = \Nsub^{-1/2} \sum_l h(\tau; t_l) \exp(- 2\pi i t_l \omega)
    \label{eqn:forward_dft_tau_omega}
\end{equation}
and its inverse
\begin{equation}
    h(\tau;t) = \Nsub^{-1/2} \sum_m h(\tau; \omega_m) \exp(2\pi i t \omega_m)
    \label{eqn:inverse_dft_omega_tau}
\end{equation}

Use these definitions to compute the partial derivative of $M$ with respect to $h^*(\tau_j;\omega_m)$
\begin{equation}
\frac{dM}{dh^*(\tau_j;\omega_m)}
 = \sum_t \frac{dM}{dh^*(\tau_j;t_l)} \frac{dh^*(\tau_j;t_l)}{dh^*(\tau_j;\omega_m)}
  = \Nsub^{-1/2} \sum_t \frac{dM}{dh^*(\tau_j;t_l)} \exp(-2\pi i t_l \omega_m)
\label{eqn:dM_dh_tau_omega}
\end{equation}

%%%%%%%%%%%%%%%%%%%%%%%%%%%%%%%%%%%%%%%%%
%%%%%%%%%%%%%%%%%%%%%%%%%%%%%%%%%%%%%%%%%

%%  Degenerate phase along the optimization trajectory

%%%%%%%%%%%%%%%%%%%%%%%%%%%%%%%%%%%%%%%%%
%%%%%%%%%%%%%%%%%%%%%%%%%%%%%%%%%%%%%%%%%

\subsection{Degenerate phase along the optimization trajectory}
\label{sec:gauge_stability}

\noindent
Section~\ref{sec:degenerate_phase} shows that $M$ is stationary with respect to $\phi_t$ at
$\phi_t=0$ (\Eqn{dM_dphi}). That result is in fact the linearization of a stronger, global
statement. Set $H'(\nu;t_l) = H(\nu;t_l)\exp(i\phi_{t_l})$, holding $H(\nu;t_{l'})$ for every
$t_{l'}\ne t_l$, $S_\mu(\alpha)$, and $g(t)$ fixed. As shown in Section~\ref{sec:degenerate_phase},
$S''_\mu(\alpha,\nu;t_l) = S'_\mu(\alpha,\nu;t_l)\, z(t_l) z^*(t_l) = S'_\mu(\alpha,\nu;t_l)$
identically, for every $\phi_{t_l}\in[0,2\pi)$ -- not merely to first order, since $z(t_l)z^*(t_l)=1$
is itself an identity, true for every $\phi_{t_l}$, not an approximation valid only for small
$\phi_{t_l}$. Because no term of \Eqn{merit_function} couples different subintegrations, $M$
depends on $t_l$ only through $S'_\mu(\cdot,\cdot;t_l)$, so $M$ is exactly invariant under this
rotation of any single subintegration's response function, at every rotation angle.

Because $h(\tau;t)$ (\Eqn{inverse_dft}) is the inverse DFT of $H(\nu;t)$ over $\nu$ at fixed $t$,
and $\phi_{t_l}$ does not depend on $\nu$, the same rotation reads
$h(\tau;t_l) \to h(\tau;t_l)\exp(i\phi_{t_l})$ for every delay $\tau$. Writing $M$ as a function of
the independent Wirtinger variables $h(\tau;t_l)$ and $h^*(\tau;t_l)$, exact invariance for every
$\phi_{t_l}$ means
\begin{equation}
M\!\left[h(\cdot;t_l)e^{i\phi_{t_l}},\, h^*(\cdot;t_l)e^{-i\phi_{t_l}}\right]
= M\!\left[h(\cdot;t_l),\, h^*(\cdot;t_l)\right]
\label{eqn:M_gauge_invariant}
\end{equation}
identically in $\phi_{t_l}$, at every $h(\cdot;t_l)$, not just at a stationary point. Differentiating
both sides with respect to $h^*(\tau_j;t_l)$, and applying the chain rule to the left-hand side
(which depends on $h^*(\tau_j;t_l)$ only through the combination $h^*(\tau_j;t_l)e^{-i\phi_{t_l}}$),
gives
\begin{equation}
\exp(-i\phi_{t_l}) \left.\frac{dM}{dh^*(\tau_j;t_l)}\right|_{h(\cdot;t_l)e^{i\phi_{t_l}}}
= \left.\frac{dM}{dh^*(\tau_j;t_l)}\right|_{h(\cdot;t_l)},
\end{equation}
i.e.
\begin{equation}
\left.\frac{dM}{dh^*(\tau_j;t_l)}\right|_{h(\cdot;t_l)e^{i\phi_{t_l}}}
= \exp(i\phi_{t_l}) \left.\frac{dM}{dh^*(\tau_j;t_l)}\right|_{h(\cdot;t_l)}
\label{eqn:grad_rotation}
\end{equation}
for every $\phi_{t_l}$: the gradient at a rotated wavefield equals the rotated gradient at the
original wavefield, exactly -- not a linearization valid only for small $\phi_{t_l}$. Expanding
\Eqn{grad_rotation} to $O(\phi_{t_l})$ and comparing coefficients reproduces \Eqn{dM_dphi};
\Eqn{grad_rotation} is the stronger, all-angle statement of the same degeneracy, verified
numerically (to floating-point precision, at every angle from $1^\circ$ to $359^\circ$ and for
$\phi_{t_l}=\pi$ exactly) against \texttt{pycyc.cyclic\_merit\_and\_grad} in
\texttt{tests/test\_fista.py}.

\Eqn{grad_rotation} holds per subintegration because the underlying invariance does: an
independent phase pattern $\{\phi_{t_l}\}$, one value per $t_l$, leaves every
$S'_\mu(\cdot,\cdot;t_l)$, and hence $M$, exactly unchanged, again because no term of
\Eqn{merit_function} couples different $t_l$. Expressed in the Doppler-delay wavefield
$h(\tau;\omega)$ (\Eqn{forward_dft_tau_omega}), such an independent pattern is \emph{not} a
pointwise operation: since $h(\tau;\omega)$ is the DFT of $h(\tau;t)$ over $t$, an independent
rotation $\{\exp(i\phi_{t_l})\}$ acts, by the convolution theorem, as a convolution of
$h(\tau;\cdot)$ with the DFT of $\{\exp(i\phi_{t_l})\}$ over $\omega$, mixing every Doppler bin.
Only a $t$-independent pattern, $\phi_{t_l}\equiv\phi$, factors out of the sum defining
$h(\tau;\omega)$ and reduces to the simple pointwise rotation $h(\tau;\omega) \to
h(\tau;\omega)\exp(i\phi)$ used in Section~\ref{sec:degenerate_phase}.

\subsubsection*{Application: the FISTA step-size estimate}

\texttt{fista.take\_fista\_step} estimates a local Lipschitz constant for the wavefield gradient
from two nearby Doppler-delay iterates, $y_n$ and $x_{n+1} = y_n - \alpha\, dM/dh^*[y_n]$, by a
discrete secant, $L \sim \|dM/dh^*[x_{n+1}] - dM/dh^*[y_n]\| / \|x_{n+1}-y_n\|$ (following the
FISTA backtracking procedure of \citealt{beck17}, chapter 10.7.1). Because \Eqn{grad_rotation}'s
degeneracy is exact, for every subintegration and every rotation angle, any component of
$x_{n+1}-y_n$ along a per-subintegration phase rotation of $y_n$ is physically unobservable and
must be removed before this ratio is evaluated. Working in the time-delay domain, where the
degeneracy is pointwise per the convolution argument above, each subintegration's best-fit
rotation
\begin{equation}
z(t_l) = \frac{\sum_\tau x^*_{n+1}(\tau;t_l)\, y_n(\tau;t_l)}
              {\left|\sum_\tau x^*_{n+1}(\tau;t_l)\, y_n(\tau;t_l)\right|}
\end{equation}
is applied to both the position difference, $z(t_l)\,x_{n+1}(\tau;t_l) - y_n(\tau;t_l)$, and, by
\Eqn{grad_rotation}, to the gradient difference,
$dM/dh^*[y_n](\tau;t_l) - z(t_l)\, dM/dh^*[x_{n+1}](\tau;t_l)$, before their norms are formed.
Applying $z(t_l)$ to only the position difference and not the gradient difference -- as an earlier
version of \texttt{take\_fista\_step} did -- leaves the gradient difference contaminated by exactly
the degenerate content the position difference had already had removed, and can inflate $L$
(and so shrink the resulting step $\alpha=1/L$) by orders of magnitude whenever an iteration
happens to be dominated by phase drift rather than genuine wavefield movement.

Because \Eqn{grad_rotation} holds at every point $h$, not only at $y_n$, it holds identically at
every iterate the FISTA momentum recursion
\begin{equation}
y_{n+1} = x_{n+1} + \frac{t_n-1}{t_{n+1}}\left(x_{n+1}-x_n\right)
\end{equation}
visits, whether or not the two iterates being compared were separated by a single gradient step
or by many. Since the same exact identity governs the gradient at every visited point, no
approximation is introduced by comparing iterates many steps apart, and no drift can accumulate
along the trajectory: confirmed to floating-point precision ($<10^{-13}$ degrees, no growth trend)
over 300 iterations of gradient descent with momentum in
\texttt{tests/test\_fista.py::test\_no\_gauge\_drift\_over\_long\_fista\_run}. This is not merely
a favorable empirical observation -- no step of \Eqn{grad_rotation}'s derivation depends on how
$h$ was reached, only on $M$ being exactly, not approximately, invariant under the rotation at
$h$ itself.

%%%%%%%%%%%%%%%%%%%%%%%%%%%%%%%%%%%%%%%%%
%%%%%%%%%%%%%%%%%%%%%%%%%%%%%%%%%%%%%%%%%

%%  Flat-response attractor at zero delay

%%%%%%%%%%%%%%%%%%%%%%%%%%%%%%%%%%%%%%%%%
%%%%%%%%%%%%%%%%%%%%%%%%%%%%%%%%%%%%%%%%%

\subsection{A flat-response attractor at zero delay}
\label{sec:flat_response_attractor}

\noindent
Sections~\ref{sec:degenerate_phase} and~\ref{sec:gauge_stability} concern directions along which
$M$ is exactly \emph{invariant}: no force, spurious or otherwise, acts along them, and
\Eqn{grad_rotation} shows the pure-phase part of that invariance cannot accumulate drift under
FISTA momentum even in principle. This section concerns a different phenomenon: a direction along
which $M$ is \emph{not} invariant, but along which the model's own bilinear structure creates a
disproportionately easy, stable minimum -- one that a natural default initial condition starts
inside of. It was diagnosed from a real-data anomaly: $h(\tau;\omega)$ develops a large,
non-physical-looking peak at $\tau=0,\,\omega=0$, together with ringing in $\omega$ that
Section~\ref{sec:gauge_stability}'s projection substantially suppresses. Unlike the ringing, the
peak itself does not decay with further iterations, and $h(0;t)$ shows why: its phase is nearly
constant and its amplitude comparable across every subintegration $t$ -- a single, essentially
$t$-independent complex number, not scintillation noise.

\subsubsection*{A $\nu$-independent response exactly reproduces the profile}

Fix a subintegration $t_l$ and suppose $H(\nu;t_l)=c$, a complex constant independent of $\nu$.
\Eqn{inverse_dft} then gives $h(\tau;t_l) = \Nchan^{1/2} c\,\delta_{\tau,0}$: a $\nu$-independent
response is, equivalently, a pure impulse at $\tau=0$. Substituting into the model,
\begin{equation}
S'_\mu(\alpha,\nu;t_l) = H(\nu+\alpha/2;t_l)H^*(\nu-\alpha/2;t_l)S_\mu(\alpha)g(t_l)
= |c|^2 S_\mu(\alpha)\, g(t_l),
\label{eqn:flat_response_model}
\end{equation}
independent of $\nu$: a flat response reproduces an exact, $\nu$-independent replica of the pulse
profile, scaled by $|c|^2 g(t_l)$, at every harmonic $\alpha$ -- it explains whatever part of
$D_\mu(\alpha,\nu;t_l)$ looks like "profile shape times a single number," with zero delay/Doppler
resolving power, using only the one real degree of freedom $|c|^2$ (the phase of $c$ is degenerate
by Section~\ref{sec:degenerate_phase} and drops out of \Eqn{flat_response_model} entirely).

\subsubsection*{The gradient at a flat response is concentrated at $\tau=0$}

Evaluate \Eqn{dSconj_dh} and \Eqn{dS_dh} at $H(\nu;t_l)=c$, so that $H(\nu\mp\alpha/2;t_l)=c$
throughout:
\begin{align}
\frac{dS_\mu^{\prime *}(\alpha,\nu;t_l)}{dh^*(\tau_j;t_l)}
&= c\,\Nchan^{-1/2}\exp(2\pi i\tau_j\nu)\exp(i\pi\tau_j\alpha)\,S_\mu^*(\alpha)\,g(t_l), \\
\frac{dS'_\mu(\alpha,\nu;t_l)}{dh^*(\tau_j;t_l)}
&= c\,\Nchan^{-1/2}\exp(2\pi i\tau_j\nu)\exp(-i\pi\tau_j\alpha)\,S_\mu(\alpha)\,g(t_l).
\end{align}
Substitute these, together with \Eqn{flat_response_model}'s residual $R_\mu(\alpha,\nu;t_l) =
|c|^2 S_\mu(\alpha) g(t_l) - D_\mu(\alpha,\nu;t_l)$, into \Eqn{dM_dh}, and use the discrete
orthogonality of the DFT basis underlying Equations~\ref{eqn:inverse_dft}--\ref{eqn:forward_dft},
\begin{equation}
\sum_\nu \exp(2\pi i \tau_j \nu) = \Nchan\,\delta_{\tau_j,0}.
\label{eqn:dft_orthogonality}
\end{equation}
Every term built from the \emph{model} piece of $R_\mu$ ($|c|^2 S_\mu(\alpha) g(t_l)$, which does
not depend on $\nu$ and so survives the $\nu$-sum only through \Eqn{dft_orthogonality}) collapses
onto $\tau_j=0$ exactly; every term built from the \emph{data} piece ($D_\mu$, which does depend on
$\nu$) does not:
\begin{align}
\left.\frac{dM}{dh^*(\tau_j;t_l)}\right|_{H=c}
= \;& 2c\,|c|^2\,g^2(t_l)\,\Nchan^{1/2}\Big(\textstyle\sum_\alpha|S_\mu(\alpha)|^2\Big)\,\delta_{\tau_j,0} \notag \\
  & - c\,g(t_l)\,\Nchan^{-1/2}\!\!\sum_{\alpha,\nu}\!\Big[S_\mu^*(\alpha)e^{i\pi\tau_j\alpha}D_\mu(\alpha,\nu;t_l)
    + S_\mu(\alpha)e^{-i\pi\tau_j\alpha}D_\mu^*(\alpha,\nu;t_l)\Big] e^{2\pi i \tau_j\nu}.
\label{eqn:flat_response_gradient}
\end{align}
The second line is present, with comparable size, at every $\tau_j$: it is simply the ordinary
data-driven pull toward matching $D_\mu$'s own delay structure, the same kind of term that acts on
$h(\tau_j;t_l)$ everywhere else along the optimization. The first line, by contrast, exists
\emph{only} at $\tau_j=0$, where discrete Fourier orthogonality forces the model's entire
self-correlation -- the profile power $\sum_\alpha|S_\mu(\alpha)|^2$, summed over all \Nchan\
channels of an exactly flat response -- to land. This is a structural property of the bilinear
shear model combined with a $\nu$-independent $H$, not a numerical coincidence, an approximation,
or a consequence of any degenerate direction: it holds even far from convergence, for any profile
$S_\mu$, any data $D_\mu$, and any subintegration.

\subsubsection*{A stable, generically non-zero equilibrium}

Setting \Eqn{flat_response_gradient} to zero at $\tau_j=0$ (where $e^{i\pi\tau_j\alpha}=e^{2\pi i
\tau_j\nu}=1$) gives, for $c\ne0$,
\begin{equation}
|c|^2 = \frac{\Nchan^{-1}\sum_{\alpha,\nu}\mathrm{Re}\left[S_\mu^*(\alpha)D_\mu(\alpha,\nu;t_l)\right]}
{g(t_l)\sum_\alpha|S_\mu(\alpha)|^2},
\label{eqn:flat_response_equilibrium}
\end{equation}
structurally the same kind of closed-form least-squares solution as the gain and profile equations
of Sections~\ref{sec:gain_variations} and~\ref{sec:intrinsic_cyclic_spectrum} -- and, because
$D_\mu$ is dominated by the pulsar's own signal, whose harmonic structure correlates with
$S_\mu(\alpha)$ essentially by construction, generically non-zero and positive. The bracketed
coefficient of $c$ in \Eqn{flat_response_gradient}'s $\tau_j=0$ term increases monotonically with
$|c|^2$ (its leading coefficient, $2g^2(t_l)\Nchan^{1/2}\sum_\alpha|S_\mu(\alpha)|^2$, is manifestly
positive), so \Eqn{flat_response_equilibrium} is the minimum of $M$ restricted to this direction,
not the unstable $c=0$ solution -- a stable fixed point that gradient descent converges to and does
not spontaneously leave. This matches the observed real-data behaviour: $|h(0;\omega{=}0)|$ grows
quickly over the first $\sim$10 iterations and then plateaus near a constant value through
convergence, rather than growing without bound the way an unstable, momentum-driven runaway would.

\subsubsection*{Why it appears at initialization, and why it persists}

\texttt{pycyc.CyclicSolver}'s default initial wavefield (\texttt{CyclicSolver.initWavefield}, before
any \texttt{first\_wavefield\_from\_best\_harmonic} refinement) sets
$h(\tau;t)=\sqrt{\Nchan\Nsub}\,\delta_{\tau,0}$ identically for every subintegration -- exactly
$H(\nu;t)=1$ for every $\nu$ and every $t$, the $H=c$ ansatz above, and, because it is exactly
$t$-independent, a pure delta in Doppler at $\omega=0$ too (\Eqn{forward_dft_tau_omega}). The first
several FISTA iterations therefore start already inside \Eqn{flat_response_equilibrium}'s basin,
for every subintegration simultaneously and at essentially the same value (since the equilibrium
depends on $t_l$ only through $D_\mu(\cdot,\cdot;t_l)$ and $g(t_l)$, which vary comparatively slowly
from one subintegration to the next) -- producing a component of $h(0;t)$ that is coherent (nearly
constant amplitude and phase) across $t$ by construction, not by any cross-subintegration coupling
in $M$ (there is none), which \Eqn{forward_dft_tau_omega} then concentrates near $\omega=0$.
(\texttt{cycsolve.py}'s actual default instead seeds the first wavefield from a best-harmonic
profile-domain estimate, \texttt{first\_wavefield\_from\_best\_harmonic=10}, a less extreme, less
purely-flat starting point that was previously found empirically to converge to a meaningfully
less delta-concentrated final wavefield than the bare-delta case -- consistent with
\Eqn{flat_response_equilibrium}, since a starting point further from the flat-response basin gives
the ordinary, non-enhanced gradient at every other $\tau_j$ correspondingly more of a head start
resolving genuine structure before $\tau=0$'s equilibrium is reached, even though
\Eqn{flat_response_equilibrium}'s own value does not depend on the initial condition at all. This
predicts that the two initializations should converge toward the same $\tau=0$ equilibrium value
give enough iterations, which has not yet been checked directly.)

Because \Eqn{flat_response_equilibrium} is a genuine local minimum of $M$ restricted to the
$\tau=0$ direction, not merely a transient the optimizer passes through, later iterations' gradient
at every \emph{other} delay -- driven only by the ordinary data term, with no analogue of the
$\delta_{\tau_j,0}$ enhancement -- has no particular incentive to relocate the power already
captured at $\tau=0$ into genuine, correctly-resolved delay-Doppler structure elsewhere: doing so
would not further reduce $M$ once \Eqn{flat_response_equilibrium} is satisfied, since the same
$\nu$-averaged, profile-correlated power is explained equally well by a concentrated impulse at
$\tau=0$ as by any physically-distributed combination of images summing to the same total. This is
why the peak plateaus rather than growing or shrinking away, and why it survives well past
convergence (in one real-data run, essentially unchanged from iteration~10 through iteration~66 --
the same run in which enabling \texttt{CS.subtract\_degenerate\_projections} independently
suppressed the surrounding ringing of Section~\ref{sec:gauge_stability} without visibly changing
this peak). The two phenomena share a location but not a cause: the ringing comes from directions
along which $M$ is exactly invariant, and momentum can in principle drift unopposed along the part
of that invariance ($\phi_t$'s delay-shift companion) that lacks \Eqn{grad_rotation}'s all-angle,
$S_\mu$-independent guarantee; this peak comes from a direction along which $M$ is \emph{not}
invariant but has a strong, stable, easily-reached minimum that happens to coincide with $\tau=0$
and with the optimizer's own initial condition.

%%%%%%%%%%%%%%%%%%%%%%%%%%%%%%%%%%%%%%%%%
%%%%%%%%%%%%%%%%%%%%%%%%%%%%%%%%%%%%%%%%%

%%  Minimal Spectral Entropy

%%%%%%%%%%%%%%%%%%%%%%%%%%%%%%%%%%%%%%%%%
%%%%%%%%%%%%%%%%%%%%%%%%%%%%%%%%%%%%%%%%%

\subsection{Minimal Spectral Entropy}

Given that the absolute phase and delay of each response function is degenerate, such that
%
$$H^\prime(\nu;t) = e^{i(\phi_t + \epsilon_t\nu)} H(\nu;t)$$ 
%
is an equally valid
solution, choose the set of phases $\{ \phi_k; 1 \le k < \Nsub\}$ and delays $\{ \epsilon_k; 1 \le k < \Nsub\}$ that concentrate power in the minimum
number of wavefield components, which is equivalent to minimizing the spectral entropy,
%
\begin{equation}
U = - \sum_j \sum_k p_{jk} \log_2 p_{jk},
\end{equation}
where $p_{jk}$ is the fractional power spectral density,
\begin{equation}
p_{jk} = P_{jk} \Sigma_P^{-1},
\end{equation}
$P_{jk} = | h^\prime(\tau_j;\omega_k) |^2$ is the power at a
given delay and Doppler shift, 
$h^\prime(\tau_j;\omega_k)$ is the forward DFT of $h^\prime(\tau;t)$,
and $h^\prime(\tau;t)$ is the inverse DFT of $H^\prime(\nu;t)$
and
%
\begin{equation}
\Sigma_P =  \sum_m \sum_n P_{mn}
\end{equation}
%
is the total power summed over all delays and Doppler shifts, which remains constant as the values of $\phi_k$ and $\epsilon_k$ are varied.
%
Minimizing the spectral entropy is aided by computing its
gradient with respect to the set of unknown phases and delays,
$\{\alpha_l\} = \{\phi_k\} \cup \{\epsilon_k\}$,
%
\begin{equation}
\frac{dU}{d\alpha_l} = - \sum_j \sum_k \frac{dp_{jk}}{d\alpha_l} \left(1+\log_2 p_{jk} \right),
\end{equation}
%
\begin{equation}
\frac{dp_{jk}}{d\alpha_l} 
= 2 \Sigma_P^{-1} \Re\left[ h^\prime(\tau_j;\omega_k) \frac{dh^{\prime*}(\tau_j;\omega_k)}{d\alpha_l} \right]
     \label{eqn:dp_dphi}
\end{equation}
%
where
%
\begin{equation} \label{eqn:dh_tau_omega_dphi}
\frac{dh^\prime(\tau_j;\omega_k)}{d\phi_l} 
  = \Nsub^{-1/2} i \exp(i\phi_l) h^\prime(\tau_j; t_l) \exp(- 2\pi i t_l \omega_k),
\end{equation}
%
\begin{equation} \label{eqn:dh_tau_omega_depsilon}
\frac{dh^\prime(\tau_j;\omega_k)}{d\epsilon_l}
  = \Nsub^{-1/2} \exp(i\phi_l) \frac{dh^\prime(\tau_j; t_l)}{d\epsilon_l}\exp(- 2\pi i t_l \omega_k),
\end{equation}
%
and
\begin{equation} \label{eqn:dh_tau_t_depsilon} 
\frac{dh^\prime(\tau_j;t_l)}{d\epsilon_l}
  = \Nchan^{-1/2} i \nu_k \exp(i\epsilon_l\nu_k) H(\nu_k; t_l) \exp(2\pi i \tau_j \nu_k),
\end{equation}
%
Substituting \Eqn{dh_tau_omega_dphi} into \Eqn{dp_dphi} and
applying $\Re[-i z] = \Im[z]$ yields
%
\begin{equation}
\frac{dp_{jk}}{d\phi_l} = 2 \Sigma_P^{-1} \Nsub^{-1/2} \Im\left[ h^\prime(\tau_j;\omega_k) h^{\prime*}(\tau_j; t_l) \exp(2\pi i t_l \omega_k) \right].
    \label{eqn:dP_dphi_expanded}
\end{equation}
%
such that
%
\begin{equation}
\frac{dU}{d\phi_l} = - 2 \Sigma_P^{-1} \sum_j \Im\left[ h^{\prime*}(\tau_j; t_l) h_w(\tau_j; t_l) \right]
\label{eqn:dU_dphi}
\end{equation}
%
where
\begin{equation}
h_w(\tau_j; t_l) = \Nsub^{-1/2} \sum_k  h^\prime(\tau_j;\omega_k)   \left(1+\log_2 p_{jk} \right)  \exp(2\pi i t_l \omega_k)
\label{eqn:h_w_def}
\end{equation}
%
is the inverse DFT of
$h^\prime(\tau_j;\omega_k) \left(1+\log_2 p_{jk} \right)$.

By the same substitution of \Eqn{dh_tau_omega_depsilon} into \Eqn{dp_dphi}
and application of $\Re[-iz]=\Im[z]$,
%
\begin{equation}
\frac{dp_{jk}}{d\epsilon_l} = 2 \Sigma_P^{-1} \Nsub^{-1/2} \Im\left[ h^\prime(\tau_j;\omega_k) \left(\frac{dh^\prime(\tau_j;t_l)}{d\epsilon_l}\right)^*_{\!/i} \exp(2\pi i t_l \omega_k) \right],
\end{equation}
%
where $\left(dh^\prime(\tau_j;t_l)/d\epsilon_l\right)_{/i}$ denotes
\Eqn{dh_tau_t_depsilon} with its explicit factor of $i$ removed, i.e.
$\Nchan^{-1/2}\nu_k\exp(i\epsilon_l\nu_k)H(\nu_k;t_l)\exp(2\pi i\tau_j\nu_k)$
-- exactly as the $\phi_l$ derivation above uses $h^\prime(\tau_j;t_l)$,
\Eqn{dh_tau_omega_dphi} with its factor of $i$ removed. Such that
%
\begin{equation}
\frac{dU}{d\epsilon_l} = - 2 \Sigma_P^{-1} \sum_j \Im\left[ \left(\frac{dh^\prime(\tau_j;t_l)}{d\epsilon_l}\right)^*_{\!/i} h_w(\tau_j; t_l) \right],
\label{eqn:dU_depsilon}
\end{equation}
%
with $h_w(\tau_j;t_l)$ the same weighted inverse DFT defined above
(\Eqn{h_w_def}): its definition depends only on the current
wavefield and $p_{jk}$, not on which of $\phi_l$ or $\epsilon_l$ is being
differentiated. \Eqn{dU_depsilon} is implemented, jointly with
\Eqn{dU_dphi} above, by
\texttt{pycyc.regularization.spectral\_entropy\_grad\_with\_delay} --
verified against \Eqn{dP_dphi_expanded}/\Eqn{dU_dphi} by requiring the two
formulations agree exactly when $\epsilon_l\equiv0$, and against both
equations directly by finite differences. It is not enabled by default:
jointly fitting $\{\phi_k\}\cup\{\epsilon_k\}$ doubles the parameter count
of the $\phi$-only fit, and at the subint counts a full observation
produces, BFGS convergence is not guaranteed (see
\texttt{CyclicSolver.minimize\_spectral\_entropy\_delay}).

%%%%%%%%%%%%%%%%%%%%%%%%%%%%%%%%%%%%%%%%%
%%%%%%%%%%%%%%%%%%%%%%%%%%%%%%%%%%%%%%%%%

%%  Degenerate-direction removal and reduced-rank jitter model

%%%%%%%%%%%%%%%%%%%%%%%%%%%%%%%%%%%%%%%%%
%%%%%%%%%%%%%%%%%%%%%%%%%%%%%%%%%%%%%%%%%

\subsection{Degenerate-direction removal and reduced-rank jitter model}
\label{sec:jitter}

Sections~\ref{sec:intrinsic_cyclic_spectrum} and~\ref{sec:jitter} address the
same underlying difficulty from two directions. Holding $S_\mu(\alpha)$
constant in $t$ is what makes \Eqn{ml_intrinsic} well posed at all: the
degenerate transform $H^\prime(\nu;t)=e^{i(\phi_t+\epsilon_t\nu)}H(\nu;t)$
leaves every individual model cyclic spectrum unchanged only if
$S_\mu(\alpha)$ is also held fixed (\S\ref{sec:degenerate_phase}), since
$\epsilon_t$ enters the model identically to a rigid shift
$S_\mu(\alpha)\rightarrow S_\mu(\alpha)e^{i\alpha\epsilon_t}$ of the pulse
profile at that one subintegration. But real pulsars exhibit genuine,
stochastic, wideband pulse-to-pulse jitter: $S_\mu(\alpha)$ does, in fact,
vary with $t$. Forcing it constant leaves that real variation nowhere to
go except into $H(\nu;t)$, where -- because jitter decorrelates between
subintegrations much faster than the interstellar scintillation timescale
-- it appears as noise that is white in $t$, and therefore white across
the entire Doppler axis $\omega$ of the reconstructed wavefield
$h(\tau,\omega)$, contaminating it well beyond the scintillation arc.

The resolution is to let $S_\mu(\alpha;t)$ vary, but constrain it to vary
only in the directions that are not degenerate with $H$. Writing the
per-subintegration profile estimate as
%
\begin{equation}
S_\mu(\alpha;t) = g(t)\,S_\mu^0(\alpha)\,e^{i\alpha\epsilon(t)} + \sum_{k=1}^{N} w_k(t)\,P_k(\alpha),
\label{eqn:jitter_model}
\end{equation}
%
where $S_\mu^0(\alpha)$ is a stable reference profile (e.g. the
subintegration- and polarization-averaged \Eqn{ml_intrinsic}), the first
term is exactly the one-dimensional degenerate direction identified above,
and the sum is a reduced-rank basis for whatever remains -- by
construction orthogonal to no particular direction a priori, but in
practice dominated by genuine jitter once the degenerate direction and the
real gain $g(t)$ (\S\ref{sec:gain_variations}) have been removed.

\paragraph{Fitting the degenerate direction.} For each subintegration,
$\epsilon(t)$ (and, when \S\ref{sec:gain_variations}'s $g(t)$ is being
modeled, a real amplitude scale) is fit by ordinary least squares against
the reference profile,
%
\begin{equation}
(\hat\epsilon(t)) = \operatorname*{arg\,min}_{\epsilon} \sum_\alpha \left| S_\mu(\alpha;t) - g(t)\,S_\mu^0(\alpha)\,e^{i\alpha\epsilon} \right|^2,
\label{eqn:fit_epsilon}
\end{equation}
%
deliberately with no free additive phase term: a rigid shift of a
\emph{real} pulse profile is exactly \Eqn{fit_epsilon}'s alpha-linear ramp
with zero intercept (the Fourier shift theorem for a real signal), so a
free phase would only fit noise and inflate the variance of the
$\hat\epsilon(t)$ estimate. \Eqn{fit_epsilon} is genuinely multi-modal once
$|\epsilon(t)|\cdot(\Nharm-1)$ exceeds $\sim\pi$ (the ramp has aliased), so
it is solved by a coarse grid search over $\epsilon$ before a BFGS
refinement, not gradient descent from a fixed starting guess -- confirmed
empirically to matter: BFGS from $\epsilon=0$ converges to the wrong local
minimum routinely at realistic $\Nharm$. This is implemented by
\texttt{pycyc.profile.fit\_profile\_shift}.

\paragraph{Reduced-rank jitter basis.} Given the residual
$R(\alpha;t) = S_\mu(\alpha;t) - g(t)\,S_\mu^0(\alpha)\,e^{i\alpha\hat\epsilon(t)}$
at every subintegration, the basis $\{P_k(\alpha)\}$ and weights
$\{w_k(t)\}$ in \Eqn{jitter_model} are the leading terms of a singular
value decomposition of $R$, whitened by the known per-harmonic noise floor
of the profile estimator. For \Eqn{ml_intrinsic}'s ordinary-least-squares
estimator, with per-channel noise variance $\sigma_D^2$ assumed independent
of radio frequency, the noise variance of a single subintegration's
$S_\mu(\alpha;t)$ is exactly
%
\begin{equation}
\mathrm{Var}\left[S_\mu(\alpha;t)\right] = \sigma_D^2 \, / \, S_\mu^{\downarrow}(\alpha;t),
\label{eqn:jitter_noise_variance}
\end{equation}
%
where $S_\mu^\downarrow(\alpha;t)$ is \Eqn{ml_intrinsic}'s (real-valued)
denominator computed for that one subintegration alone -- the standard
result for ordinary least squares against a known design coefficient with
i.i.d.\ per-sample noise. Whitening $R(\alpha;t)$ by
$\sqrt{\mathrm{Var}[S_\mu(\alpha;t)]}$ and comparing the eigenvalues of the
resulting sample covariance against the Marchenko--Pastur edge,
$(1+\sqrt{\Nharm/\Nsub})^2$, gives an objective choice of $N$: under the
null hypothesis that $R$ is pure noise at the variance given by
\Eqn{jitter_noise_variance}, no eigenvalue should exceed this bound, so
any that do indicate genuine structure rather than noise. This is
implemented by \texttt{pycyc.jitter.fit\_jitter\_basis}.

\paragraph{Outer loop.} Because $S_\mu^0(\alpha)$, $\epsilon(t)$, $g(t)$,
and the jitter basis all depend on the current estimate of $H(\nu;t)$
-- and, early in an optimization, on how much of the residual is genuine
jitter versus how much is simply what is still missing from an
unconverged $H$ -- \Eqn{jitter_model} is refreshed periodically rather
than fit once: \texttt{CyclicSolver.outer\_loop} alternates the wavefield
fit with a profile-pooling step and a refresh of
\Eqn{fit_epsilon}--\Eqn{jitter_noise_variance}, deferring the jitter model
entirely for an initial number of warm-up passes so it is not fit against
transient wavefield-model error. The retained rank $N$ and the principal
angle between successive passes' bases (a standard subspace-distance
metric, via QR followed by the singular values of the two orthonormalized
bases' inner product) serve as the natural convergence diagnostic: both
are expected to stabilize once the outer loop has converged.












\bibliography{journals,modrefs,psrrefs,crossrefs}
\bibliographystyle{aasjournal}

\end{document}


\subsection{Minimal Spectral Entropy}

Given that the absolute phase of each response function is degenerate, such that 
$$h^\prime(\tau;t) = \exp(i\phi_t)h(\tau;t)$$ 
is an equally valid
solution, choose the set of phases $\{ \phi_k; 1 \le k < \Nsub\}$ that concentrate power in the minimum
number of wavefield components, which is equivalent to minimizing the spectral entropy,
%
\begin{equation}
U = - \sum_j \sum_k p_{jk} \log_2 p_{jk},
\end{equation}
where $p_{jk}$ is the fractional power spectral density,
\begin{equation}
p_{jk} = P_{jk} \Sigma_P^{-1},
\end{equation}
$P_{jk} = | h^\prime(\tau_j;\omega_k) |^2$ is the power at a
given delay and Doppler shift, 
$h^\prime(\tau_j;\omega_k)$ is the forward DFT of $h^\prime(\tau;t)$,
and
%
\begin{equation}
\Sigma_P =  \sum_m \sum_n P_{mn}
\end{equation}
%
is the total power summed over all delays and Doppler shifts, which remains constant
as the values of $\phi_k$ are varied.
%
Minimizing the spectral entropy is aided by computing its
gradient with respect to the unknown phases
%
\begin{equation}
\frac{dU}{d\phi_l} = - \sum_j \sum_k \frac{dp_{jk}}{d\phi_l} \left(1+\log_2 p_{jk} \right),
\end{equation}
%
where
%
\begin{equation}
\frac{dp_{jk}}{d\phi_l} 
= 2 \Sigma_P^{-1} \Re\left[ h^\prime(\tau_j;\omega_k) \frac{dh^{\prime*}(\tau_j;\omega_k)}{d\phi_l} \right]
     \label{eqn:dp_dphi}
\end{equation}
%
and
%
\begin{equation}
    \frac{dh^\prime(\tau_j;\omega_k)}{d\phi_l} = \Nsub^{-1/2} i \exp(i\phi_l) h(\tau_j; t_l) \exp(- 2\pi i t_l \omega_k)
    \label{eqn:dh_tau_omega_dphi}.
\end{equation}
%
Substituting \Eqn{dh_tau_omega_dphi} into \Eqn{dp_dphi} and
applying $\Re[-i z] = \Im[z]$ yields
%
\begin{equation}
\frac{dp_{jk}}{d\phi_l} = 2 \Sigma_P^{-1} \Nsub^{-1/2} \Im\left[ h^\prime(\tau_j;\omega_k) h^{\prime*}(\tau_j; t_l) \exp(2\pi i t_l \omega_k) \right].
    \label{eqn:dP_dphi_expanded}
\end{equation}
%
such that
%
\begin{equation}
\frac{dU}{d\phi_l} = - 2 \Sigma_P^{-1} \sum_j \Im\left[ h^{\prime*}(\tau_j; t_l) h_w(\tau_j; t_l) \right]
\end{equation}
%
where
\begin{equation}
h_w(\tau_j; t_l) = \Nsub^{-1/2} \sum_k  h^\prime(\tau_j;\omega_k)   \left(1+\log_2 p_{jk} \right)  \exp(2\pi i t_l \omega_k) 
\end{equation}
%
is the inverse DFT of
$h^\prime(\tau_j;\omega_k) \left(1+\log_2 p_{jk} \right)$.


\bibliography{journals,modrefs,psrrefs,crossrefs}
\bibliographystyle{aasjournal}

\end{document}


\subsection{Gradient of Maximum Likelihood Intrinsic Profile}

This way leads to madness!  Adding the second terms of equations \ref{eqn:dSconj_dh_degenerate}
and \ref{eqn:dS_dh_degenerate} to the definition of the merit function causes the problem to become
degenerate, as described in Section 2.1 of \cite{wdv13}.  If you want to try it out anyhow, express \Eqn{ml_intrinsic} as
%
\begin{equation}
    S_\mu(\alpha) =\frac{ S_\mu^\uparrow(\alpha) }{ S_\mu^\downarrow(\alpha) },
\label{eqn:ml_profile}
\end{equation}
%
where the (complex-valued) numerator,
%
\begin{equation}
    S_\mu^\uparrow(\alpha) = \sum_{t,\nu} D(\alpha,\nu;t) H^*(\nu+\alpha/2;t) H(\nu-\alpha/2;t) g(t)
\end{equation}
%
and the (real-valued) denominator,
%
\begin{equation}
    S_\mu^\downarrow(\alpha) = \sum_{t,\nu} | H(\nu+\alpha/2;t) |^2 | H(\nu-\alpha/2;t) |^2 g^2(t).
\end{equation}
%

Noting that $S_\mu^\downarrow(\alpha)$ is real-valued,
\begin{equation}
\frac{dS_\mu^*(\alpha)}{dh^*(\tau_j;t_l)}
    = \frac{ 1 }
                              { S_\mu^\downarrow(\alpha) }\frac{dS_\mu^{\uparrow *}(\alpha)}{dh^*(\tau_j;t_l)}
    - \frac{ S_\mu^{\uparrow *}(\alpha) }
                              { (S_\mu^\downarrow(\alpha))^2 }\frac{dS_\mu^\downarrow(\alpha)}{dh^*(\tau_j;t_l)},
\label{eqn:dSxconj_dh}
\end{equation}
%
where
%
\begin{equation}
\frac{dS_\mu^{\uparrow *}(\alpha)}{dh^*(\tau_j;t_l)}
= \left[ \Nchan^{-1/2} \sum_\nu D^*(\alpha,\nu;t_l) H(\nu+\alpha/2;t_l) \exp(2\pi i \tau_j \nu) \right] \exp(-\pi i \tau_j \alpha) g(t_l)
\end{equation}
%
(the bit in square brackets is an inverse DFT) and
%
\begin{equation}
\begin{split}
\frac{dS_\mu^\downarrow(\alpha)}{dh^*(\tau_j;t_l)}
&= g^2(t_l) \left[ \Nchan^{-1/2} \sum_\nu |H(\nu+\alpha/2;t_l)|^2 H(\nu-\alpha/2;t_l) \exp(2\pi i \tau_j \nu) \right] \exp(-\pi i \tau_j \alpha) \\
&+ g^2(t_l) \left[ \Nchan^{-1/2} \sum_\nu |H(\nu-\alpha/2;t_l)|^2 H(\nu+\alpha/2;t_l) \exp(2\pi i \tau_j \nu) \right] \exp(\pi i \tau_j \alpha)    
\end{split}
\end{equation}
%
(the bits in square brackets are inverse DFTs).
%
Substitution of the second term of \Eqn{dSconj_dh} into the first term of \Eqn{dM_dh} yields
\begin{equation}
g(t_l) \sum_\alpha \frac{dS_\mu^*(\alpha)}{dh^*(\tau_j;t_l)}
\sum_\nu H(\nu-\alpha/2;t_l) H^*(\nu+\alpha/2;t_l) R(\alpha,\nu;t_l).
\label{eqn:dM_dh_Sxconj_term}
\end{equation}
%
Similarly,
%
\begin{equation}
\frac{dS_\mu(\alpha)}{dh^*(\tau_j;t_l)}
    = \frac{ 1 }{ S_\mu^\downarrow(\alpha) }\frac{dS_\mu^\uparrow(\alpha)}{dh^*(\tau_j;t_l)}
    - \frac{ S_\mu^\uparrow(\alpha) }
            { (S_\mu^\downarrow(\alpha))^2 }\frac{dS_\mu^\downarrow(\alpha)}{dh^*(\tau_j;t_l)}
\end{equation}
%
where
%
\begin{equation}
\frac{dS_\mu^\uparrow(\alpha)}{dh^*(\tau_j;t_l)}
= \left[ \Nchan^{-1/2} \sum_\nu D(\alpha,\nu;t_l) H(\nu-\alpha/2;t_l) \exp(2\pi i \tau_j \nu) \right] \exp(\pi i \tau_j \alpha) g(t_l).
\end{equation}
%
Substitution of the second term of \Eqn{dS_dh} into the second term of \Eqn{dM_dh} yields
\begin{equation}
g(t_l) \sum_\alpha \frac{dS_\mu(\alpha)}{dh^*(\tau_j;t_l)}
\sum_\nu H^*(\nu-\alpha/2;t_l) H(\nu+\alpha/2;t_l) R^*(\alpha,\nu;t_l).
\label{eqn:dM_dh_Sx_term}
\end{equation}
%
Note that the sum over $\nu$ in \Eqn{dM_dh_Sx_term} is simply the complex conjugate of the sum over $\nu$ in
\Eqn{dM_dh_Sxconj_term}.

\bibliography{journals,modrefs,psrrefs,crossrefs}
\bibliographystyle{aasjournal}

%% This command is needed to show the entire author+affiliation list when
%% the collaboration and author truncation commands are used.  It has to
%% go at the end of the manuscript.
%\allauthors

%% Include this line if you are using the \added, \replaced, \deleted
%% commands to see a summary list of all changes at the end of the article.
%\listofchanges

\end{document}

% End of file `sample631.tex'.
